% Auto-generated from Chapter-46-Refugees-Hear-of-Aisens-Death.md
% Source: C:\Users\PCGAME\Desktop\story&universes\chain-weapon-story\finished-manuscript\manuscriptbaseforchapters\Chapter-46-Refugees-Hear-of-Aisens-Death.md

\chapter{Refugees Hear of Aisen's Death}
\renewcommand{\CurrentChapterNum}{46}
\POV{}{Emma / Refugees}



The bread was wrong.

Emma knew it by touch before she knew it by taste—the crust yielding a fraction too early beneath her thumb, the interior holding its warmth unevenly, the grain's structure compressed in a way that spoke of flour milled too coarse and water added too fast. She had baked bread for twenty-three years. Her hands knew the language of dough the way a musician's hands knew the language of strings: by pressure, by resistance, by the specific conversation between the material and the maker that produced, when both participants were honest, something that could sustain a body.

This bread was not honest. But it was bread.

She broke it and distributed the pieces to the children gathered at the long table in the refectory tent—a word too grand for the structure, which was canvas stretched over a frame of green pine poles, the canvas patched in three places with sailcloth and the poles still weeping sap that caught the lamplight in amber beads. The children received the bread with the automatic hands of people who had learned not to examine food before eating it, because examination was a luxury that belonged to a life before, a life where bread came from ovens with names—\emph{Emma's, on the square}—and the bread had a provenance and the provenance was a form of trust and the trust was a form of home.

There was no home now. There was the camp.

The camp was a place the way a scar was a place—defined not by what it was but by what had happened to produce it. It occupied a stretch of floodplain above the Miren delta, the ground flat and damp and given to a mineral smell that rose from the clay when the rain came, which was often. The tents numbered in the hundreds. The population shifted—families arriving, families dispersing to relatives or work or the grey anonymity of cities that absorbed refugees the way rivers absorbed tributaries, by dilution. Emma had been here for fourteen months. She was, by the camp's accelerated calendar, an elder.

She was also, by the camp's informal taxonomy, a baker. Not officially—officially she was Displaced Person 4,417, Thessaly District, Female, Category B (Able-bodied, Skilled Trade, No Dependents Above Age of Labour)—but functionally. She baked. She had three ovens, built by her own hands from the clay of the floodplain and fired with the same technique her grandmother had taught her grandmother: a slow cure, the heat rising in controlled increments over three days, the clay learning its shape through the gradual education of fire. The ovens were not her mother's ovens. They were not the stone-floored, three-chambered ovens of the bakery on the square in Thessaly, the village, not the province. But they baked bread, and the bread fed people, and the feeding was the thing that she could do, the single act of competence that the upheaval had not taken from her.

The news arrived on a Tuesday.

Not dramatically. Not carried by a rider or announced by a herald or delivered in the formal language of institutional communication. It came the way all news came to the margins: sideways, through the channels of the displaced, the informal network of travellers and traders and the simply restless who moved between camps and settlements carrying information the way wind carried seeds—without intention, without agenda, by the mechanics of movement through a populated landscape.

A woman named Darya, who sold cloth remnants from a cart she had salvaged from a failed textile workshop in Greenvale, mentioned it at the water pump.

"The quiet one," she said. She was filling a bucket, her hands red from the cold, her voice carrying the flat declarative tone of someone reporting weather. "The one with the chain. They say he fell at the western ridge."

Emma's hands stopped. She was rinsing flour from her forearms, the water cold enough to ache, the flour dissolving into a thin paste that the current carried away. Her hands stopped and the water continued and the paste continued dissolving and the dissolving was indifferent to the stopping, because the world did not pause when a fact arranged itself inside a person's understanding and changed the architecture of what that person knew.

"When?" she asked.

"A week. Maybe more. News travels slow."

A week. Maybe more. The quiet boy who had taught strangers to survive had been dead for a week, maybe more, and the world had continued baking and washing and filling buckets and dissolving flour, and the continuation was not cruelty but reality, the planet's stubborn insistence on rotating regardless of who had stopped rotating with it.

She carried the bucket back to the ovens. The dough waiting beneath the damp cloth had risen in the faithful, indifferent way that dough rose when flour and water and yeast were given enough warmth to continue being themselves, and the faithfulness of it angered her for a breath and then steadied her, because bread did not ask whether the hands that shaped it had earned tomorrow, and neither, she thought, had the boy with the chain. When the smallest child at the table asked, later, who he had been, she gave him the smallest true answer first. "Someone who noticed," she said, dividing the loaf while the oven breathed heat against her scarred wrists. "Someone who taught other people to notice in time." The children ate. The fact entered them with the bread.